\section{Lecture III: Orbital integrals}\label{s:orb}
The significance of affine Springer fibers in representation theory is demonstrated by their close relationship with orbital integrals. Orbital integrals are certain integrals that appear in the harmonic analysis of $p$-adic groups. Just as conjugacy classes of a finite group are fundamental to understanding its representations, orbital integrals are fundamental to understanding representations of $p$-adic groups. In certain cases, orbital integrals can be interpreted as counting points on affine Springer fibers.


\subsection{Integration on a $p$-adic group}\label{ss:int}

\subsubsection{The setup} Let $F$ be a local non-archimedean field, i.e., $F$ is either a finite extension of $\QQ_{p}$ or a finite extension of $\FF_{p}\lr{t}$. Then $F$ has a discrete valuation $\val: F^{\times}\to \ZZ$ which we normalize to be surjective. Let $\calO_{F}$ be the valuation ring of $F$ and $k$ be the residue field. Therefore, unlike in the previous sections,  $k$ is a finite field. We assume that $\chark$ is large with respect to the groups in question.



\subsubsection{Haar measure and integration} Let $G$ be an algebraic group over $F$. The topological group $G(F)$ is locally compact and totally disconnected. It has a right invariant Haar measure $\mu_{G}$ which is unique up to a scalar. For a measurable subset $S\subset G(F)$, we denote its volume under $\mu_{G}$ by $\vol(S,\mu_{G})$. Fixing a compact open subgroup $K_{0}\subset G(F)$, we may normalize the Haar measure $\mu_{G}$ so that $K_{0}$ has volume $1$. For example, if we choose an integral model $\calG$ of $G$ over $\calO_{F}$, we may take $K_{0}=\calG(\calO_{F})$.  

With the Haar measure $\mu_{G}$ one can integrate smooth (i.e., locally constant) compactly supported functions on $G(F)$ with complex values. We denote this function space by $\sS(G(F))$ (where $\sS$ stands for Schwarz). For $f\in\sS(G(F))$, the integral
\begin{equation*}
\int_{G(F)}f\mu_{G}
\end{equation*} 
can be calculated as follows. One can find a subgroup $K\subset K_{0}$ of finite index such that $f$ is right $K$-invariant, i.e., $f(gx)=f(g)$ for all $g\in G(F)$ and $x\in K$ (see Exercise \ref{exIII:K inv}). Then the integral above becomes a weighted counting in the coset $G(F)/K$:
\begin{equation*}
\int_{G(F)}f\mu_{G}=\vol(K,\mu_{G})\sum_{[g]\in G(F)/K}f(g)=\frac{1}{[K_{0}:K]}\sum_{[g]\in G(F)/K}f(g).
\end{equation*}
It is easy to check that the right side above  is independent of the choice of $K$ as long as $f$ is right $K$-invariant and $K$ has finite index in $K_{0}$.

\subsubsection{Variant} Let $H\subset G$ be a subgroup defined over $F$ together with a Haar measure $\mu_{H}$ on it. Consider a function $f\in\sS(H(F)\bs G(F))$, i.e., $f$ is a left $H(F)$-invariant, locally constant function on $G(F)$ whose support is compact modulo $H(F)$, we may define the integral
\begin{equation}\label{int quot}
\int_{H(F)\bs G(F)}f\frac{\mu_{G}}{\mu_{H}}.
\end{equation} 
This integral is calculated in the following way. Again we choose a finite index subgroup $K\subset G(F)$ such that $f$ is right $K$-invariant. Then the integral \eqref{int quot} can be written as a weighted sum over double cosets $H(F)\bs G(F)/K$:
\begin{equation*}
\int_{H(F)\bs G(F)}f\frac{\mu_{G}}{\mu_{H}}=\frac{1}{[K_{0}:K]}\sum_{[g]\in H(F)\bs G(F)/K}\frac{f(g)}{\vol(H(F)\cap gKg^{-1}, \mu_{H})}.
\end{equation*}

\subsection{Orbital integrals} 

\subsubsection{Definition of orbital integrals} We continue with the setup of \S\ref{ss:int}. We denote the Lie algebra of $G$ by $\frg(F)$ to emphasize that it is a vector space over $F$. Let $\varphi\in\sS(\frg(F))$ and $\gamma\in\frg(F)$. Consider the map $G(F)\to \frg(F)$ given by $g\mapsto \Ad(g^{-1})\gamma$. Then the composition $f: g\mapsto \varphi(\Ad(g^{-1})\gamma)$ is a smooth function on $G(F)$.  Then $f$ is locally constant, left invariant under the centralizer $G_{\gamma}(F)$ of $\gamma$ in $G(F)$, and has compact support modulo $G_{\gamma}(F)$. 

Fix Haar measures $\mu_{G}$ on $G(F)$ and $\mu_{G_{\gamma}}$ on $G_{\gamma}(F)$. The following integral is then a special case of \eqref{int quot} (except that we write the integration variable $g\in G_{\gamma}(F)\bs G(F)$ explicit below while not in \eqref{int quot})
\begin{equation*}
O_{\gamma}(\varphi):=\int_{G_{\gamma}(F)\bs G(F)}\varphi(\Ad(g^{-1})\gamma)\frac{\mu_{G}}{\mu_{G_{\gamma}}}.
\end{equation*}
Such integrals are called {\em orbital integrals} on the Lie algebra $\frg(F)$. We may similarly define orbital integrals on the group $G(F)$ by replacing $\gamma$ with an element in $G(F)$ and $\varphi$ with an element in $\sS(G(F))$.

\subsubsection{Specific situation} For the rest of the section we will restrict to the following situation. Let $G$ be a split reductive group over $F$. We may fix an integral model of $G$ by base changing the corresponding Chevalley group scheme from $\ZZ$ to $\calO_{F}$. In the following we will regard $G$ as a reductive group scheme over $\calO_{F}$. We normalize the Haar measure $\mu_{G}$ on $G(F)$ by requiring that $K_{0}=G(\calO_{F})$ have volume $1$.

The Lie algebra $\frg(F)$ contains a canonical lattice $\frg(\calO_{F})$ coming from the integral model over $\calO_{F}$. We will be most interested in the orbital integral of the characteristic function $\varphi=\one_{\frg(\calO_{F})}$ of the lattice $\frg(\calO_{F})$. 

\subsubsection{The centralizer of $\gamma$} Suppose $\gamma$ is regular semisimple so that its centralizer $G_{\gamma}$ is a torus over $F$. Let $F^{\ur}=F\wh{\otimes}_{k}\kbar$, which is a complete discrete valuation field whose residue field is algebraically closed. We continue to let $t$ denote a uniformizer of $F$, which is also a uniformizer in $F^{\ur}$. Using $t$, the construction in \S\ref{sss:L gamma} gives an embedding $\xcoch(G_{\gamma}\otimes_{F}F^{\ur})\incl  G_{\gamma}(F^{\ur})$ whose image we still denote by $\L_{\gamma}$. This embedding being $\Gk$-equivariant, $\L_{\gamma}$ carries an action of $\Gk$. We think of $\L_{\gamma}$ as an \'etale group scheme over $k$, then the notation $\L_{\gamma}(k)$ makes sense, and it is just the $\Gk$-invariants in $\L_{\gamma}$ if we regard the latter as a plain group. Then $\L_{\gamma}(k)\subset G_{\gamma}(F)$ is a discrete and cocompact subgroup.


\subsubsection{Centralizers in $\GL_{n}$}\label{sss:GLn orb} Let $G=\GL_{n}$ and let $\gamma$ be a regular semisimple element in $\frg(F)$ which is not necessarily diagonalizable over $F$. Assume either $\textup{char}(F)=0$ or $\chark>n$. As in \S\ref{sss:GLn centralizer}, the characteristic polynomial $P(x)=x^{n}+a_{1}x^{n-1}+\cdots+a_{n}\in F[x]$ of $\gamma$ is separable, hence the $F$-algebra $F[x]/(P(x))$ is isomorphic to a product of fields $F_{1}\times \cdots \times F_{m}$. We have
\begin{equation*}
G_{\gamma}(F)\cong F^{\times}_{1}\times\cdots\times F^{\times}_{m}.
\end{equation*}
In this case, the lattice $\L_{\gamma}(k)\subset G_{\gamma}(F)=F^{\times}_{1}\times\cdots\times F^{\times}_{m}$ consists of elements of the form $(t^{d_{1}},\cdots, t^{d_{m}})$ for $(d_{1},\cdots, d_{m})\in\ZZ^{m}$. The quotient $\L_{\gamma}(k)\bs G_{\gamma}(F)$ is isomorphic to $\prod_{i=1}^{m}F^{\times}_{i}/t^{\ZZ}$. Each $F^{\times}_{i}/t^{\ZZ}$ fits into an exact sequence
\begin{equation}\label{Fi seq}
0\to \calO^{\times}_{F_{i}}\to F^{\times}_{i}/t^{\ZZ}\to \ZZ/e_{i}\ZZ\to 0
\end{equation}
where $e_{i}$ is the ramification degree of the extension $F_{i}/F$, therefore the quotient $\L_{\gamma}(k)\bs G_{\gamma}(F)$ is compact.

\subsubsection{Orbital integrals in terms of counting}\label{sss:X gamma} Consider the following subset of $G(F)/G(\calO_{F})$
\begin{equation*}
X_{\gamma}:=\{[g]\in G(F)/G(\calO_{F})|\Ad(g^{-1})\gamma\in\frg(\calO_{F})\}.
\end{equation*}
This is a set-theoretic version of the affine Springer fiber. 

The group $G_{\gamma}(F)$ acts on $X_{\gamma}$ by the rule $G_{\gamma}(F)\ni h:[g]\mapsto[hg]$. For any free abelian group $L\subset G_{\gamma}(F)$, its action on $G(F)/G(\calO_{F})$ by left translation is free (because the stabilizers are necessarily finite), hence it acts freely on $X_{\gamma}$.

More generally, for any discrete cocompact subgroup $L\subset G_{\gamma}(F)$, the quotient groupoid $L\bs X_{\gamma}$ is {\em finitary}, i.e.,  it has finitely many isomorphism classes and the automorphism group of each object is finite. For a finitary groupoid $Y$, we define the cardinality of $Y$ to be 
\begin{equation}\label{card groupoid}
\#Y:=\sum_{y\in \textup{Ob}(Y)/\cong}\frac{1}{\#\Aut(y)}
\end{equation}


The next lemma follows directly from the definitions, whose proof is left to the reader as Exercise \ref{exIII:Ogamma}.
\begin{lemma}\label{l:Ogamma} Let $\gamma$ be a regular semisimple element in $\frg(F)$. Let $L\subset G_{\gamma}(F)$ be any discrete cocompact subgroup. We have
\begin{equation*}
O_{\gamma}(\one_{\frg(\calO_{F})})=\frac{1}{\vol(G_{\gamma}(F)/L,\mu_{G_{\gamma}})}\#\left(L\bs X_{\gamma}\right)
\end{equation*}
with the cardinality on the right side interpreted as in \eqref{card groupoid}.
\end{lemma}





\subsubsection{The case $G=\GL_{n}$ and fractional ideals}\label{sss:GLn orb formula} We continue with the situation in \S\ref{sss:GLn orb}. Under the identification of $G(F)/G(\calO_{F})$ with the set of $\calO_{F}$-lattices in $F^{n}$ (see \eqref{Gr lattices}), we have
\begin{equation*}
X_{\gamma}\cong\{\textup{lattices }\L\subset F^{n}|\gamma \L\subset \L\}.
\end{equation*} 
The bijection sends $[g]\in X_{\gamma}$ to the lattice $\L=g\calO^{n}_{F}$. 

We give another interpretation of $X_{\gamma}$. Let $P(x)=x^{n}+a_{1}x^{n-1}+\cdots+a_{n}$ be the characteristic polynomial of $\gamma$. Let
\begin{equation*}
A=\calO_{F}[x]/(P(x))
\end{equation*}
be the commutative $\calO_{F}$-subalgebra of $\frg(F)$ generated by $\gamma$. The ring of total fractions $\Frac(A)$ is a finite \'etale $F$-algebra of degree $n$, and $A$ is an order in it. The canonical action of $A$ on $F^{n}$ realizes $F^{n}$ as a free $\Frac(A)$-module of rank $1$.  Recall a {\em fractional $A$-ideal} is a finitely generated $A$-submodule $M\subset \Frac(A)$. If we choose an element $e\in F^{n}$ as a basis for the $\Frac(A)$-module structure, we get a bijection
\begin{equation}\label{frac ideals}
\{\mbox{fractional $A$-ideals}\}\bij X_{\gamma}
\end{equation}
which sends $M\subset \Frac(A)$ to the $\calO_{F}$-lattice $M\cdot e\subset F^{n}$. 



Using the algebra $A$, we have a canonical isomorphism
\begin{equation*}
G_{\gamma}(F)\cong \Frac(A)^{\times}.
\end{equation*}
This isomorphism intertwines the action of $G_{\gamma}(F)$ on $X_{\gamma}$ by left translation and the action of $\Frac(A)^{\times}$ on the set of fractional $A$-ideals by multiplication. 

When $A$ happens to be a product of Dedekind domains (i.e., $A$ is the maximal order in $\Frac(A)$), all fractional $A$-ideals are principal, which is the same as saying that the action of $G_{\gamma}(F)\cong \Frac(A)^{\times}$ on $X_{\gamma}$ is transitive. In general, principal fractional ideals form a homogeneous space $\Frac(A)^{\times}/A^{\times}$ under $G_{\gamma}(F)$; the difficulty in counting $X_{\gamma}$ in general is caused by the singularity of the ring $A$.


We normalize the Haar measure $\mu_{G_{\gamma}}$ on $G_{\gamma}(F)\cong\prod_{i}F^{\times}_{i}$ so that $\prod_{i}\calO^{\times}_{F_{i}}$ gets volume $1$. Choose a uniformizer $t_{i}\in F_{i}$, we may form the lattice $L_{0}=t^{\ZZ}_{1}\times \cdots\times t^{\ZZ}_{m}\subset G_{\gamma}(F)$. Now $G_{\gamma}(F)/L_{0}\cong \prod_{i}\calO_{F_{i}}^{\times}$ has volume $1$. Therefore Lemma \ref{l:Ogamma} gives
\begin{equation}\label{GLn L0}
O_{\gamma}(\one_{\frg(\calO_{F})})=\#(L_{0}\bs X_{\gamma}).
\end{equation}
Using \eqref{frac ideals}, we may interpret \eqref{GLn L0} as saying that $O_{\gamma}(\one_{\frg(\calO_{F})})$ is the number of fractional $A$-ideals up to multiplication by the powers of the $t_{i}$'s. 


